\section{Model Kontrol Akses dan Relevansi ABAC pada IoT}

Model kontrol akses digunakan untuk menentukan bagaimana suatu sistem mengatur hak akses pengguna atau perangkat terhadap sumber daya tertentu. Salah satu model paling awal adalah \textit{Discretionary Access Control} (DAC), yaitu model yang memberikan kewenangan kepada pemilik sumber daya untuk memberikan atau menolak izin akses kepada identitas atau kelompok tertentu melalui mekanisme seperti \textit{Access Control List} (ACL) \parencite{diaz_authorization_2025}. Model ini relatif fleksibel karena pengelolaan izin dilakukan langsung oleh pemilik objek, namun pendekatan tersebut menyulitkan penerapan kebijakan keamanan yang terpusat dan konsisten pada lingkungan berskala besar. \textcite{mao_zero-trust_2025} menyatakan bahwa DAC menawarkan fleksibilitas tetapi terlalu sulit untuk menerapkan kebijakan keamanan yang terpadu, sehingga dapat meenyebabkan \textit{permission creep} dan inkonsistensi kebijakan keamanan.

Sementara itu, \textit{Mandatory Access Control} (MAC) menggunakan pendekatan berbasis label atau level keamanan yang ditentukan secara terpusat oleh administrator sistem. Dalam model ini, hak akses ditentukan berdasarkan perbandingan atribut keamanan pengguna dan objek terhadap level otorisasi yang telah ditetapkan \parencite{diaz_authorization_2025}. MAC menawarkan tingkat keamanan yang tinggi karena kebijakan akses dikendalikan secara ketat oleh sistem, tetapi model ini bersifat \textit{rigid} dan kurang adaptif terhadap perubahan kebutuhan akses yang dinamis. \textcite{mao_zero-trust_2025} menjelaskan bahwa konfigurasi kebijakan MAC terlalu \textit{rigid} dan sulit menyesuaikan diri dengan kebutuhan akses yang sering berubah pada lingkungan \textit{cloud} yang dinamis dan kompleks. Kondisi tersebut membuat MAC kurang sesuai diterapkan pada lingkungan IoT yang memiliki karakteristik heterogen, dinamis, dan membutuhkan fleksibilitas tinggi dalam pengelolaan akses. \textcite{ragothaman_access_2023} juga menekankan bahwa lingkungan IoT memiliki kebutuhan akses kontrol yang bersifat dinamis, heterogen, dan membutuhkan granularitas tinggi dalam pengambilan keputusan akses.

Selain DAC dan MAC, \textit{Role-Based Access Control} (RBAC) dikembangkan untuk menyederhanakan pengelolaan hak akses melalui konsep peran (\textit{role}). Dalam RBAC, pengguna tidak diberikan izin secara langsung, melainkan memperoleh izin akses melalui peran yang dimiliki. \textcite{ferraiolo_proposed_2001} menjelaskan bahwa konsep dasar RBAC adalah pengguna dikelompokkan berdasarkan peran, izin diberikan kepada peran, dan pengguna memperoleh izin dengan menjadi anggota suatu peran. Pendekatan ini mempermudah administrasi hak akses karena pengelolaan izin akses dapat dilakukan pada level peran, bukan pada setiap pengguna secara individual. RBAC juga banyak digunakan pada sistem berskala besar karena mampu mendukung pengelolaan hak akses yang lebih terstruktur dan efisien \parencite{ferraiolo_proposed_2001}. Namun, pada lingkungan IoT yang sangat dinamis dan heterogen, RBAC memiliki keterbatasan dalam menangani kebutuhan akses yang bersifat kontekstual. \textcite{diaz_authorization_2025} menyatakan bahwa RBAC dapat mengalami kesulitan pada lingkungan IoT karena sering membutuhkan tambahan faktor seperti kondisi lingkungan, waktu, dan lokasi dalam proses pengambilan keputusan akses. Selain itu, \textcite{ameer_hybrid_2023} menegaskan bahwa pada lingkungan \textit{cloud} dan IoT yang kompleks, RBAC dapat mengalami masalah \textit{role explosion} ketika jumlah peran dan kombinasi izin akses meningkat secara signifikan.

Untuk mengatasi keterbatasan model akses tradisional yang hanya melakukan evaluasi sebelum akses diberikan, dikembangkan konsep \textit{Usage Control} (UCON) yang mendukung kontrol akses yang lebih dinamis dan berkelanjutan. \textcite{park_uconabc_2004} menjelaskan bahwa UCON merupakan generalisasi dari kontrol akses yang mencakup otorisasi, obligasi, kondisi, kontinuitas, dan mutabilitas. Berbeda dengan model tradisional yang umumnya hanya melakukan keputusan akses pada saat permintaan awal, UCON juga mendukung \textit{ongoing control} selama sesi penggunaan berlangsung. Dengan demikian, hak akses dapat berubah secara dinamis berdasarkan perubahan atribut, kondisi lingkungan, maupun perilaku pengguna selama sesi aktif. UCON juga mendukung pembaruan atribut secara terus-menerus sebagai konsekuensi dari aktivitas akses yang sedang berlangsung \parencite{park_uconabc_2004}. Meskipun menawarkan fleksibilitas dan kontrol yang lebih kaya dibandingkan DAC, MAC, dan RBAC, implementasi UCON cenderung lebih kompleks karena membutuhkan pemantauan atribut, kondisi, dan evaluasi kebijakan secara kontinu selama proses penggunaan sumber daya berlangsung.

\textit{Attribute-Based Access Control} (ABAC) menjadi salah satu pendekatan yang paling relevan untuk IoT karena keputusan akses dibuat berdasarkan atribut subjek, objek, aksi, dan lingkungan yang berkaitan dengan permintaan akses \parencite{hu_guide_2014}. Dengan pendekatan ini, keputusan akses tidak hanya ditentukan oleh identitas atau peran, tetapi juga oleh konteks seperti lokasi, waktu, jenis perangkat, status keamanan perangkat, dan jenis aksi yang diminta. \textcite{diaz_authorization_2025} menjelaskan bahwa ABAC mampu menyediakan \textit{fine-grained access control} karena izin dapat ditentukan secara lebih spesifik berdasarkan kombinasi atribut yang relevan. Selain itu, \textcite{kalaria_adaptive_2024} menunjukkan bahwa kontrol akses IoT perlu bersifat adaptif dan \textit{context-aware} agar mampu menyesuaikan kebijakan akses terhadap perubahan kondisi lingkungan secara dinamis. Karakteristik tersebut membuat ABAC lebih sesuai untuk IoT dibandingkan DAC, MAC, maupun RBAC yang cenderung lebih statis.

Meskipun demikian, penerapan ABAC pada IoT juga memiliki tantangan tersendiri. ABAC membutuhkan banyak atribut dari berbagai sumber untuk mengevaluasi setiap permintaan akses, sehingga proses pengambilan atribut dari PIP, \textit{cloud}, atau penyedia atribut lain dapat meningkatkan \textit{overhead} komunikasi dan latensi. \textcite{cremonezi_improving_2022} menunjukkan bahwa \textit{attribute retrieval} pada ABAC di lingkungan \textit{fog-IoT} dapat menjadi mahal pada skenario IoT yang besar dan \textit{mobile}, sehingga strategi \textit{caching} atribut diperlukan untuk meningkatkan performa. Selain itu, kualitas keputusan akses sangat bergantung pada \textit{freshness} atribut yang digunakan. \textcite{wang_edge-enabled_2025} menegaskan bahwa keputusan akses pada \textit{edge-IAM} dapat menjadi tidak reliabel apabila sistem mengasumsikan atribut selalu mutakhir meskipun kondisi aktual telah berubah. Oleh karena itu, meskipun ABAC paling sesuai untuk kebutuhan kontrol akses IoT yang dinamis, penerapannya tetap membutuhkan strategi pengelolaan atribut yang efisien agar performa dan keamanan sistem dapat terjaga secara seimbang. 